\section{本章小结}
\label{sec:feas-summary}

本章的问题是怎样让每个迭代点满足原约束，以及保持或恢复可行性要花多少计算。预算例给出了同一问题上的对比，圆周例说明了凸集之外的情形。

投影梯度把梯度步拉回最近点，残差$\bm{G}_\alpha$为零等价于变分不等式，$L$光滑时目标差有$1/T$的界，强凸时距离线性收缩。条件梯度只用线性最小化，靠凸组合保持可行，产生顶点的短组合，间隙$\Gamma$在凸问题上上界目标差。镜像下降把欧氏距离换成Bregman距离，在单纯形上给出保持正性的指数更新；一般次梯度分析针对平均点，只有相对光滑时目标才逐步下降。

GRG面向光滑非线性约束，通过可逆的基矩阵局部消元，按约化梯度移动非基本变量，再解方程恢复基本变量，碰界或病态时换基。它依赖的是局部可解性，与前三种方法的闭凸集假设不同，两边的保证不能互相借用。

各种方法都要分别检查原约束误差和合适的最优性残差。内层投影、线性最小化或恢复不精确时，误差要进入递推和停止检验。可行迭代还需要一个可行的起点，这个起点本身有时也要专门构造。
